\chapter{\ifenglish Background Knowledge and Theory\else ทฤษฎีที่เกี่ยวข้อง\fi}

\section{Cryptography}
Cryptography หรือ วิทยาการรหัสลับ เป็นศาสตร์ที่ว่าด้วยการแปลงข้อความปกติให้กลายเป็นข้อความลับ โดยข้อความลับเป็นข้อความที่ผู้อื่นนอกเหนือจากคู่สนทนาที่ต้องการไม่สามารถเข้าใจได้
โดยทั่วไป กระบวนการปกป้องข้อมูลจะเริ่มต้นก่อนการจัดเก็บหรือการส่งข้อมูล โดยขั้นตอนนี้เรียกว่า การเข้ารหัส (Encryption) 
ซึ่งเป็นการนำข้อมูลต้นฉบับมาผ่านกระบวนการทางคณิตศาสตร์ร่วมกับกุญแจ (Key) ที่เป็นค่าตัวเลขสุ่ม 
เพื่อแปลงเป็นข้อมูลที่ถูกเข้ารหัส (Ciphertext) ในทางกลับกัน เมื่อต้องการอ่านข้อมูลดังกล่าว 
จะต้องนำข้อมูลที่ถูกเข้ารหัสมาผ่านกระบวนการทางคณิตศาสตร์ร่วมกับกุญแจอีกครั้ง เพื่อแปลงกลับเป็นข้อมูลดั้งเดิม 
ซึ่งขั้นตอนนี้เรียกว่า การถอดรหัส (Decryption)

\subsection{Symmetric and Asymmetric Encryption}
การเข้ารหัสในสมัยใหม่ถูกแบ่งออกเป็น 2 ประเภทหลักตามวิธีการใช้กุญแจดังนี้ คือ การเข้ารหัสแบบสมมาตร (Symmetric Encryption) และ การเข้ารหัสแบบอสมมาตร (Asymmetric Encryption)

\subsubsection{Symmetric Encryption}
เป็นกระบวนการเข้ารหัสที่ใช้กุญแจ (Key) ดอกเดียวกันทั้งในขั้นตอนการเข้ารหัสและการถอดรหัส ซึ่งเป็นวิธีที่มีประ-สิทธิภาพสูงและประมวลผลได้อย่างรวดเร็ว ยากที่จะทำการถอดรหัสหากไม่ทราบกุญแจ 
แต่มีข้อจำกัดที่สำคัญคือปัญหาในการกระจายกุญแจ 
เนื่องจากผู้ส่งและผู้รับจำเป็นต้องตกลงและส่งมอบกุญแจผ่านช่องทางที่ปลอด-ภัยล่วงหน้า

\subsubsection{Asymmetric Encryption}
เป็นกระบวนการเข้ารหัสที่ใช้กุญแจสองดอกที่มีความสัมพันธ์กันทางคณิตศาสตร์ที่ยากที่จะแก้ไขแต่ง่ายที่จะตรวจสอบได้
ได้แก่ กุญแจสาธารณะ (Public Key) สำหรับการเข้ารหัส และกุญแจส่วนตัว (Private Key) สำหรับการถอดรหัส 
วิธีการนี้ถูกคิดค้นขึ้นเพื่อแก้ปัญหาการกระจายกุญแจของการเข้ารหัสแบบสมมาตร 
และได้รับการยอมรับอย่างแพร่หลายในปัจจุบัน รวมถึงการประยุกต์ใช้ในระบบการโหวตออนไลน์
แต่ด้วยการที่ต้องใช้กุญแจที่มีขนาดใหญ่มากในการเข้ารหัสเพื่อที่จะทำให้การถอดรหัสทำได้ยาก
ทำให้การประมวลผลช้ากว่าการเข้ารหัสแบบสมมาตร


\subsection{Computational Security and Information-Theoretic Security}

\subsubsection{Computational Security}
Computational Security เป็นความปลอดภัยที่อาศัยสมมติฐานที่ว่าปัญหาทางคณิตศาสตร์บางประเภทมีความยากและซับซ้อนเกินกว่าที่จะแก้ไขได้ภายในระยะเวลาและทรัพยากรที่เหมาะสม 
แนวคิดนี้เป็นรากฐานสำคัญของการเข้ารหัสสมัยใหม่ส่วนใหญ่ ไม่ว่าจะเป็น RSA, ECC, AES 
ตลอดจนอัลกอริทึมการเข้ารหัสยุคหลังควอนตัม (Post-Quantum Cryptography: PQC) ระดับความปลอดภัยของระบบเหล่านี้จะถูกประเมินจากเวลาและทรัพยากรทางคอมพิวเตอร์ที่จำเป็นต้องใช้ในการเจาะระบบ 
อย่างไรก็ตาม เมื่อเทคโนโลยีคอมพิวเตอร์ควอนตัมมีความก้าวหน้ามากขึ้น ระบบการเข้ารหัสเหล่านี้จึงเผชิญกับความเสี่ยงที่จะถูกเจาะทำลาย 
แม้ว่าการเข้ารหัสแบบ PQC จะถูกออกแบบมาเพื่อลดความเสี่ยงดังกล่าว แต่วิธีการเหล่านั้นก็ยังคงพึ่งพาสมมติฐานเชิงคำนวณเป็นหลัก จึงไม่สามารถป้องกันการโจมตีจากระบบที่มีพลังการประมวลผลอย่างไม่จำกัดได้


\subsubsection{Information-Theoretic Security}
Information-Theoretic Security เป็นระดับความปลอดภัยที่ไม่ได้พึ่งพาข้อสมมติฐานด้านขีดความสามารถในการคำนวณใดๆ 
โดยระบบสามารถรับประกันได้ว่าผู้โจมตีจะไม่สามารถถอดรหัสข้อความได้เลย แม้ว่าจะมีพลังการประมวลผลทางคอมพิวเตอร์สูงอย่างไม่จำกัดก็ตาม 
หลักการสำคัญของความปลอดภัยเแบบนี้คือการที่ทำให้รูปแบบการเข้ารหัสยังคงปลอดภัยจากการโจมตีและยังคงรับประกันว่าระดับความปลอดภัยจะไม่เสื่อมถอยลงตามกาลเวลา 
แม้เทคโนโลยีจะพัฒนาไปมากเพียงใดในอนาคต อย่างไรก็ตาม ความท้าทายหลักของระบบนี้คือความซับซ้อนใ่นการนำไปใช้งานจริง
เนื่องจากผู้ใช้งานจำเป็นต้องปฏิบัติตามกฎของโปรโตคอลอย่างเคร่งครัด เพื่อรักษาความสมบูรณ์ของระบบโดยไม่ลดทอนระดับความปลอดภัยลง

\subsection{One Time Pad}



\section{Quantum Key Distribution}
\subsection{Quantum Property}
\subsubsection{Superposition}
Superposition เป็นสมบัติของควอนตัมที่ทำให้สถานะของระบบควอนตัมสามารถอยู่ในสถานะต่างๆ ได้พร้อมกัน
จนกว่าจะมีการวัดค่าเกิดขึ้น คุณสมบัตินี้เป็นกลไกสำคัญที่ทำให้ผู้ส่งสามารถเข้ารหัสข้อมูลลงในอนุภาคโฟตอน (Photon) ได้
\subsubsection{No Cloning Theorem}
หลักการพื้นฐานทางกลศาสตร์ควอนตัม หากไม่รู้ว่าอนุภาคควอนตัมอยู่ในสถานะใด จะไม่สามารถสร้างสำเนาของอนุภาคนั้นได้อย่างสมบูรณ์แบบ
ซึ่งทำให้ผู้ที่ดักฟังข้อมูลไม่สามารถคัดลอกข้อมูลไปได้
\subsubsection{Measurement}
การวัดค่าสถานะทางควอนตัมจะทำให้สถานะของระบบควอนตัมเปลี่ยนแปลงไปสู่สถานะใดสถานะหนึ่ง หรือก็คือการทำให้ Superposition สลายไป
ซึ่งผลลัพธ์ที่ได้จากการวัดค่าสถานะทางควอนตัมจะขึ้นอยู่กับรูปแบบวิธีที่ใช้ในการวัดค่า

\subsection{The BB84 Protocol}

\section{\ifenglish%
\ifcpe CPE \else ISNE \fi knowledge used, applied, or integrated in this project
\else%
ความรู้ตามหลักสูตรซึ่งถูกนำมาใช้หรือบูรณาการในโครงงาน
\fi
}

อธิบายถึงความรู้ และแนวทางการนำความรู้ต่างๆ ที่ได้เรียนตามหลักสูตร ซึ่งถูกนำมาใช้ในโครงงาน

\section{\ifenglish%
Extracurricular knowledge used, applied, or integrated in this project
\else%
ความรู้นอกหลักสูตรซึ่งถูกนำมาใช้หรือบูรณาการในโครงงาน
\fi
}

อธิบายถึงความรู้ต่างๆ ที่เรียนรู้ด้วยตนเอง และแนวทางการนำความรู้เหล่านั้นมาใช้ในโครงงาน
